\subsection{Definable flips for balls}\label{sec:definable-flips-for-balls}

% lax begin Lax185574Proofs.VCDuality.vc_duality
The following lemma is a corollary of the $(p,q)$-Theorem for bounded VC-dimension from~\cite{Matousek:p-q-theorem}.
See~\cite[Corollary~6.13]{simon_book} or~\cite[full version, Theorem~E.2]{flip-width} for self-contained proofs.
% lax end

For a graph $G$ and sets $S,X\subseteq V(G)$, we say that $S$ \emph{dominates} $X$ if every vertex of $X \setminus S$ has a neighbor in $S$, and that $S$ \emph{antidominates} $X$ if every vertex of $X \setminus S$ has a nonneighbor in $S$.

% lax begin Lax185574.VCDuality
\begin{lemma}[Bounded-VC duality]\label{lem:vc-duality}
    For every $d\in\N$ there is some $k\in\N$ with the following property.
    Let $G$ be a graph of VC-dimension at most $d$, and let $X,Y\subseteq V(G)$ be two (possibly overlapping) sets with at least one of $X,Y$ finite.
    Then there exists a set $\sample(X,Y)$ of size at most $k$ such that at least one of the following holds:
    \begin{enumerate}
        \item $\sample(X,Y)$ is contained in $Y$ and dominates $X$;
        \item $\sample(X,Y)$ is  contained in $X$ and antidominates $Y$.
    \end{enumerate}
\end{lemma}
% lax end
We remark that in the sources we have given, the above lemma is stated either for finite set systems or finite binary relations.
The translation to finite graphs is straightforward, by either viewing the graph as a set system of neighborhoods or working with the adjacency relation.
Moreover, the statement lifts to infinite graphs if at least one of the sets $X$ and $Y$ is finite, as the other one can be replaced by its finite subset representing all neighborhoods in the  finite one.
(However, the statement may fail if both $X$ and $Y$ are infinite, as witnessed by the bi-infinite half-graph with vertices $\mathbb{Z} \times \{1,2\}$ and edges between $(x,1)$ and $(y,2)$ if $x\leq y$.)

It was observed already in~\cite{DBLP:conf/lics/BonnetDGKMST22} that \Cref{lem:vc-duality} is useful to obtain definable flips.
We will use it in a similar way to deduce the remaining result of \Cref{sec:closures-main}.

\lemBallFlip*

Note that, given the assumptions of the lemma, it is easy to find a (non-definable) $(2^t+t+1)$-flip $H'$ of $G$ in which $\Ball_{r+1}^{H'}(C) \cap A = \emptyset$: group the vertices of $A$ and of $\Ball_r^G(C)$ by their neighborhoods in the other set, which gives at most $2^t$ and $t$ groups, respectively, and flip every group of $A$ with every group of $\Ball_r^G(C)$ that it is complete to.
So the challenge is really to get such a flip in a definable way.

We also remark that \Cref{lem:ball-flip} can be seen as a finitization of Lemma~7.3 in the full version of~\cite{flippers}.
That lemma works in a more model-theoretic setting, where $A$ is a (possibly infinite) model instead of a set.
It uses a version of \Cref{lem:vc-duality} which also works for infinite $X,Y$, under the stronger
assumption that $G$ excludes a half-graph of order $d$.

% lax begin Lax185574Proofs.BallFlip.exists_ball_flip
\begin{proof}[Proof of \Cref{lem:ball-flip}]
    Put $B\coloneqq \Ball_r^G(C)$.
    Let $\AA$ be the partition of $A$ that groups vertices by their neighborhood in $B$, and let $\BB$ be the partition of $B$ that groups vertices by their neighborhood in $A$.
    By assumption $|\BB|\leq t$.
    The neighborhood in $B$ of a vertex of $A$ is a union of parts of $\BB$, so $|\AA|\leq2^t$.
    Moreover, every pair $P\in\AA$, $Q\in\BB$ is homogeneously connected: is either complete or anticomplete.

    Let $k \leq O_d(1)$ be given by \Cref{lem:vc-duality}.
    Pick any set $S$ which contains
    \begin{itemize}
        \item an arbitrary representative $\rep(P)\in P$ for each part $P\in\AA\cup\BB$, and
        \item
        the set $\sample(X,Y)$ supplied by \Cref{lem:vc-duality}, for every $X\in\AA$ and $Y\in\AA\cup\{V(G)\setminus B\}$
        (the lemma is applicable because $X \subseteq A$ is finite).
    \end{itemize}
    We can pick such a set with $|S|\le z$ for some $z\le O_{k,t}(1)\le O_{d,t}(1)$.
    We will construct $H$ as an $S$-flip of $G$; in particular, $H$ is $z$-definable.

    Let $\PP_S$ be the partition of $V(G)$ into atomic types over $S$.
    Write $x\sim y$ when $x$ and $y$ are in the same $\PP_S$-part.
    Observe the following two properties of $\PP_S$.

    \begin{claim}
    $\PP_S$ refines $\AA$ on $A$ and $\BB$ on $B$.
    \end{claim}
    \begin{claimproof}
        If two vertices $a$ and $a'$ of $A$ are in different $\AA$-parts, then their adjacencies differ on some vertex $b \in B$.
        Let $P$ be the $\BB$-part containing $b$.
        The vertex $\rep(P) \in S$ has the same adjacencies to $A$ as $b$ and hence distinguishes $a$ and $a'$ in $\PP_S$.
        Thus, $\PP_S$ refines $\AA$ on $A$ and symmetrically also $\BB$ on $B$.
    \end{claimproof}

    \begin{claim}\label{clm:dagger}
    Let $X\in \AA$ and $Y\in \AA\cup\set{V(G)-B}$.
    There do not exist vertices $x,y,x',y'$ with:
    \begin{equation}\tag{$\dagger$}\label{eq:ball-flip-crossing}
        x\in X,\quad y\in Y,\quad x\sim x',\quad y\sim y',
        \quad Y\subseteq N^{\overline G}(x'),\quad X\subseteq N^G(y').
    \end{equation}
    \end{claim}
    \begin{claimproof}
    Suppose that $\sample(X,Y)$ lies in $Y$ and dominates $X$.
    Then there is some  $s\in\sample(X,Y)\cap Y$ with $sx\in E(G)$.
    If
    $Y\subseteq N^{\overline G}(x')$,
then $sx'\notin E(G)$.
Thus, $x\nsim x'$.

The case when $\sample(X,Y)$ lies in $X$ and antidominates $Y$, is symmetric, by swapping $X$ with $Y$ and $G$ with $\overline G$.
    Both cases are illustrated in \Cref{fig:ball-flip-crossing}.
    \end{claimproof}
    \begin{center}
        \definecolor{bfAnti}{RGB}{200,40,40}%
\definecolor{bfHl}{RGB}{0,90,200}%
{\setlength{\fboxsep}{6pt}\makebox[\linewidth]{\colorbox{white}{%
\begin{tikzpicture}[
        vertex/.style={circle, fill=black, inner sep=0pt, minimum size=3.2pt},
        blob/.style={draw=black!70, line width=0.6pt, fill=black, fill opacity=0.04},
        complete/.style={draw=black, line width=0.6pt},
        anticomplete/.style={draw=bfAnti, dash pattern=on 2.5pt off 2pt, line width=0.6pt,
            shorten <=2pt},
        sameType/.style={black!60, densely dotted, line width=0.7pt},
        lbl/.style={inner sep=1pt},
        sampleSet/.style={draw=bfHl, line width=0.6pt, fill=bfHl, fill opacity=0.12},
        hlEdge/.style={draw=bfHl, line width=1.2pt},
        hlNonedge/.style={draw=bfHl, dash pattern=on 3pt off 2pt, line width=1.2pt},
    ]
    \def\bx{1.7}
    \def\rx{1.1}
    \def\ry{0.8}
    \def\off{0.7}
    \def\h{2.2}
    \def\panelsep{6.3}

    \pgfmathsetmacro{\bfunit}{\linewidth / (2*\panelsep + 2*(\bx+\rx))}
    \begin{scope}[x=\bfunit pt, y=\bfunit pt]

    \pgfmathsetmacro{\tphi}{atan2(\h/\ry, (2*\bx+\off)/\rx)}
    \pgfmathsetmacro{\tdelta}{acos(1/sqrt(((2*\bx+\off)/\rx)^2 + (\h/\ry)^2))}

    \def\bfSkipComplete{-1}
    \def\bfSkipAnti{-1}
    \def\bfSimLabel{$\sim$}

    \def\bfCrossingBase{%
        \coordinate (X) at (-\bx, 0);
        \coordinate (Y) at ( \bx, 0);
        \coordinate (x) at (-\bx-\off, 0.15);
        \coordinate (y) at ( \bx+\off, 0.15);
        \coordinate (x') at (-\bx-\off, \h);
        \coordinate (y') at ( \bx+\off, \h);

        \pgfmathsetmacro{\px}{2*\bx+\off}
        \pgfmathsetmacro{\fanfrom}{atan2(\ry*sin(\tphi+\tdelta) - \h, \rx*cos(\tphi+\tdelta) - \px)}
        \pgfmathsetmacro{\fanto}{atan2(\ry*sin(\tphi-\tdelta) - \h, \rx*cos(\tphi-\tdelta) - \px)}
        \foreach \i in {0,...,5} {
            \pgfmathsetmacro{\fanang}{\fanfrom + \i*(\fanto - \fanfrom)/5}
            \pgfmathsetmacro{\fanA}{(cos(\fanang))^2/\rx^2 + (sin(\fanang))^2/\ry^2}
            \pgfmathsetmacro{\fanB}{2*(\px*cos(\fanang)/\rx^2 + \h*sin(\fanang)/\ry^2)}
            \pgfmathsetmacro{\fanC}{\px^2/\rx^2 + \h^2/\ry^2 - 1}
            \pgfmathsetmacro{\fanlen}{(-(\fanB) - sqrt(max(0, \fanB*\fanB - 4*\fanA*\fanC)))/(2*\fanA)}
            \ifnum\i=\bfSkipComplete\relax\else
                \draw[complete] (y') -- ++(\fanang:\fanlen);
            \fi
            \ifnum\i=\bfSkipAnti\relax\else
                \draw[anticomplete, dash phase={mod(\i,2)*2.25pt}] (x') -- ++(180-\fanang:\fanlen);
            \fi
        }

        \draw[blob] (X) ellipse [x radius=\rx, y radius=\ry];
        \draw[blob] (Y) ellipse [x radius=\rx, y radius=\ry];
        \node[lbl, below=2pt] at ($(X) + (0, -\ry)$) {$X$};
        \node[lbl, below=2pt] at ($(Y) + (0, -\ry)$) {$Y$};

        \draw[sameType] (x') -- node[lbl, left=1pt, black] {\bfSimLabel} (x);
        \draw[sameType] (y') -- node[lbl, right=1pt, black] {\bfSimLabel} (y);

        \foreach \v in {x, y, x', y'}
            \node[vertex] at (\v) {};
        \node[lbl, below=2pt] at (x) {$x$};
        \node[lbl, below=2pt] at (y) {$y$};
        \node[lbl, above=2pt] at (x') {$x'$};
        \node[lbl, above=2pt] at (y') {$y'$};
    }

    \begin{scope}[name prefix=a-]
        \bfCrossingBase
    \end{scope}

    \begin{scope}[shift={(\panelsep, 0)}, name prefix=b-]
        \def\bfSkipAnti{3}
        \def\bfSimLabel{}
        \bfCrossingBase
        \pgfmathsetmacro{\bfsang}{180 - (\fanfrom + \bfSkipAnti*(\fanto - \fanfrom)/5)}
        \pgfmathsetmacro{\bfslen}{(2*\bx + \off - 0.4)/cos(\bfsang)}
        \coordinate (s) at ($(x') + (\bfsang:\bfslen)$);
        \draw[sampleSet] (s) ellipse [x radius=0.45, y radius=0.3];
        \draw[hlEdge] (s) -- (x);
        \draw[hlNonedge] (s) -- (x');
        \foreach \v in {x, x'}
            \node[vertex] at (\v) {};
        \node[vertex, fill=bfHl] at (s) {};
    \end{scope}

    \begin{scope}[shift={(2*\panelsep, 0)}, name prefix=c-]
        \def\bfSimLabel{}
        \bfCrossingBase
        \pgfmathsetmacro{\bfsang}{\fanfrom + 3*(\fanto - \fanfrom)/5}
        \pgfmathsetmacro{\bfslen}{-(2*\bx + \off - 0.4)/cos(\bfsang)}
        \coordinate (s) at ($(y') + (\bfsang:\bfslen)$);
        \draw[sampleSet] (s) ellipse [x radius=0.45, y radius=0.3];
        \draw[hlNonedge] (s) -- (y);
        \draw[hlEdge] (s) -- (y');
        \foreach \v in {y, y'}
            \node[vertex] at (\v) {};
        \node[vertex, fill=bfHl] at (s) {};
    \end{scope}

    \coordinate (bbtop) at (current bounding box.north);
    \coordinate (bbbot) at (current bounding box.south);
    \pgfresetboundingbox
    \path[use as bounding box]
        (-\bx-\rx,0 |- bbbot) rectangle (2*\panelsep+\bx+\rx,0 |- bbtop);
    \end{scope}
\end{tikzpicture}%
}}}%

        \captionof{figure}{%
            Left: the forbidden configuration~\eqref{eq:ball-flip-crossing}.
            Middle: the case of the proof of \Cref{clm:dagger} where $\sample(X,Y)$ (marked in blue) lies in $Y$.
            Right: the case where $\sample(X,Y)$ lies in $X$.}
        \label{fig:ball-flip-crossing}
    \end{center}

    Define the $S$-flip $H$ of $G$ where a pair $(P,Q)$ of parts of $\PP_S$ is flipped if and only if there is an edge in $G$ with endpoints in $P$ and in $Q$ and also in $A$ and in $B$. More formally, $H[P,Q]=G[P,Q]$ whenever
    $E(G) \cap PQ \cap AB =
    \emptyset,$
    and otherwise, $H[P,Q]=\overline G[P,Q]$.

    The following three properties will yield the conclusion of the lemma:
    \begin{enumerate}[label=(P\arabic*), ref=P\arabic*]
        \item\label{itm:ball-flip-P1} $A$ and $\Ball_{r}^G(C) $ are nonadjacent in $H$.
        \hfill (proven in \Cref{clm:ball-flip-P1})
        \item\label{itm:ball-flip-P2} $\Ball_{r-1}^H(C)=\Ball_{r-1}^G(C)$.
        \hfill (proven in \Cref{clm:ball-flip-P2})
        \item\label{itm:ball-flip-P3}$\Ball_{r-1}^G(C)$ and $V(G) \setminus \Ball_{r}^G(C)$ are nonadjacent in $H$.
        \hfill (proven in \Cref{clm:ball-flip-P3})
    \end{enumerate}
    If $r = 0$, then already \ref{itm:ball-flip-P1} alone yields the conclusion. This justifies the implicit assumption $r \geq 1$ made in  \ref{itm:ball-flip-P2} and \ref{itm:ball-flip-P3}.
    By combining \ref{itm:ball-flip-P2} and \ref{itm:ball-flip-P3}, we derive:
    \begin{enumerate}[label=(P\arabic*), ref=P\arabic*,start = 4]
        \item\label{itm:ball-flip-P4}  $\Ball_{r}^H(C) \subseteq \Ball_{r}^G(C)$.
    \end{enumerate}
    The conclusion of the lemma thus follows:
    \[
       A \cap \Ball_{r+1}^H(C)
       \
    =
    \
       A \,\cap\, \Ball_1^H\big(\Ball_r^H(C)\big)
       \  \underset{\rule{0pt}{8pt}\text{\ref{itm:ball-flip-P4}}}{\subseteq} \
       A \cap \Ball_1^H\big(\Ball_r^G(C)\big)
       \  \underset{\rule{0pt}{8pt}\text{\ref{itm:ball-flip-P1}}}{=\vphantom{\subseteq}} \
       A \cap \Ball_r^G(C)
       \  \underset{\rule{0pt}{8pt}\text{ass.}}{=\vphantom{\subseteq}} \
       \emptyset.
    \]
    It only remains to establish \ref{itm:ball-flip-P1}--\ref{itm:ball-flip-P3}.

    \begin{claim}[\ref{itm:ball-flip-P1}]\label{clm:ball-flip-P1}
        $A$ and $B$ are nonadjacent in $H$.
    \end{claim}
    \begin{claimproof}
By contrapositive, suppose $ab\in E(H)$ for some $a \in A$ and $b \in B$.
    By definition of $H$,  $ab\in E(G)$ implies $ab\notin E(H)$, so
     $ab\notin E(G)$.
    Since $ab\in E(H)$, the parts of $\PP_S$ containing $a$ and $b$ are flipped, that is, there is some $a'b'\in E(G)$ with $a' \in A$ and $b' \in B$,
    and
    \begin{center}
        $(a\sim a'$ and $b\sim b')$
        \qquad
        or
        \qquad
        $(a\sim b'$ and $b\sim a')$.
    \end{center}
    The first case cannot happen
    as $\PP_S$ refines $\AA$ on $A$ and $\BB$ on $B$, but the two sets $\{a,a'\}$ and $\{b,b'\}$ are inhomogeneous as $ab\notin E(G)$ and $a'b'\in E(G)$.

    Thus, assume $a\sim b'$ and $b\sim a'$.
    Let $X,Y\in\AA$ contain $a',a$, respectively.
    Then $a'b'\in E(G)$ implies that $b'$ is complete to $X$, whereas  $ab\notin E(G)$ implies that $b$ is anticomplete to $Y$ (both by homogeneity). This contradicts \eqref{eq:ball-flip-crossing} with $(x,y,x',y')=(a',a,b,b')$.
\end{claimproof}

    For brevity, we define $D\coloneqq \Ball_{r-1}^G(C)$. Recall that $\PP_S(v)$ is the part of $\PP_S$ containing $v$.

    \begin{claim}\label{clm:disjoint-from-d}
        If $b$ is adjacent to some $a\in A$ in $G$ then
         $\PP_S(b)\cap D=\emptyset$.
    \end{claim}
    \begin{claimproof}
            The vertices in $B \cap \PP_S(b)$ have the same neighborhood on $A$ as~$b$. In particular, they are all adjacent to $a$, whereas $D \subseteq B$ is anticomplete to $A$ since $\Ball^G_1(D)=\Ball^G_r(C)$ is disjoint from $A$ by assumption.
    \end{claimproof}

    \begin{claim}[\ref{itm:ball-flip-P2}]\label{clm:ball-flip-P2}
        $\Ball_{r-1}^H(C)=D$.
    \end{claim}
    \begin{claimproof}
    The claim is trivially true for $r = 1$, hence assume $r\geq 2$.
    Consider any  $ab\in E(G)$ with $a \in A$ and $b \in B$ witnessing a flipped pair of $\PP_S$-parts $\PP_S(a)$ and $\PP_S(b)$.

    \begin{enumerate}
    \item \emph{$\PP_S(b)$ is disjoint from $\Ball_{r-2}^G(C)$:}

    This is implied by \Cref{clm:disjoint-from-d}, because $\Ball_{r-2}^G(C) \subseteq \Ball_{r-1}^G(C) =D$.

    \item \emph{$\PP_S(a)$ is disjoint from $\Ball_{r-2}^G(C)$:}

    Let $Q\in\BB$ contain $b$.
    Every vertex of $Q$ is adjacent to $a$ and hence lies in $B\setminus D$.
    Thus, $\rep(Q)\in S$ is adjacent to $a$ and nonadjacent to every vertex of $\Ball_{r-2}^G(C)$.
    Hence, $\rep(Q)$ witnesses that $a$ is in a different $\PP_S$-part than all the vertices in $\Ball_{r-2}^G(C)$.
    \end{enumerate}

    This means that whenever two parts $\PP_S(a)$ and $\PP_S(b)$ are flipped, then neither of them intersects~$\Ball_{r-2}^G(C)$.
    Hence, the vertices in $\Ball_{r-2}^G(C)$ are incident to the same edges in $H$ as in $G$, which gives the desired conclusion
    $\Ball_{r-1}^H(C) = \Ball_{r-1}^G(C)$.
    \end{claimproof}

    For brevity, we define $\overline B\coloneqq V(G)\setminus B$.

    \begin{claim}[\ref{itm:ball-flip-P3}]\label{clm:ball-flip-P3}
        $D$ and $\overline B$ are nonadjacent in $H$.
    \end{claim}

    \begin{claimproof}
            Otherwise, an edge $dz\in E(H)$ with $d \in D$ and $z \in \overline B$ is witnessed by some $ab\in E(G)$ with $a \in A$ and $b \in B$ (because $D$ and $\overline B$ are nonadjacent in $G$).
    By \Cref{clm:disjoint-from-d}, the $\PP_S$-part of $b$ is disjoint from $D$, so necessarily $d\sim a$ and $z\sim b$.
    Let $X\in\AA$ contain $a$.
    The vertex $b$ is complete to $X$, whereas $d$ is anticomplete to $\overline B$, contradicting \eqref{eq:ball-flip-crossing} for $(X,\overline B)$ with $(x,y,x',y')=(a,z,d,b)$.
    \end{claimproof}
    As described above, establishing \ref{itm:ball-flip-P1}--\ref{itm:ball-flip-P3} finishes the proof.
\end{proof}
% lax end
